\documentclass[a4paper,11pt]{article}

% --- Packages ---
\usepackage[utf8]{inputenc}
\usepackage[T1]{fontenc}
\usepackage{geometry}
\usepackage{graphicx}
\usepackage{svg}
\usepackage{booktabs}
\usepackage{tabularx}
\usepackage{siunitx}
\usepackage{amsmath}
\usepackage{amssymb}
\usepackage{hyperref}
\usepackage{caption}
\usepackage{fancyhdr}
\usepackage{subcaption}
\usepackage{float}

% --- Configuration ---
\geometry{top=2.5cm, bottom=2.5cm, left=2.5cm, right=2.5cm}

\sisetup{
    output-decimal-marker = {,},
    group-digits = true,
    per-mode = symbol
}

% SVG Settings
% DO NOT use \svgsetup{convert=true} with the standard 'svg' package.
% Just ensure Inkscape is in PATH and compile with --shell-escape.
% The 'svg' package will automatically call Inkscape to convert .svg to .pdf.

\hypersetup{
    colorlinks=true,
    linkcolor=blue,
    filecolor=magenta,      
    urlcolor=cyan,
    pdftitle={5-DOF Robotic Arm Documentation},
    pdfauthor={Benjamin McEwan}
}

% --- DOCUMENT -------------------------------------------------------------------------------------
\begin{document}

\title{\textbf{5-DOF Robotic Arm}}
\author{Benjamin McEwan}
\date{\today}
\maketitle

\tableofcontents


% --- 1: INTRODUCTION ------------------------------------------------------------------------------
\newpage
\section{Introduction}
\begin{figure}[h!]
    \centering
    % NO 'crop' option here. Just the filename.
    % Ensure you cropped the SVG in Inkscape first!
    \includesvg[width=0.5\linewidth]{svg/arm.svg} 
    \caption{Full arm assembly}
    \label{fig:arm}
\end{figure}

This project is a simple yet modular and easily modifiable 5-DOF robotic arm intended to be a platform for the teaching and learning of basic robotics control systems. The project uses five MG996R servo motors as they are a common hobbyist robotics component. Figure \ref{fig:arm} shows the final assembled arm in a resting position.

The arm does not possess an actuated gripper as it has been designed with each servo placed directly within the arm's joints in order to avoid the complexities of gearboxes or belt drives. This allows for the stated focus on control systems and means that any issues encountered in actuation of the arm can be immediately attributed to the servos themselves, streamlining the physical troubleshooting process.


% --- 2: DRAWINGS ----------------------------------------------------------------------------------
\newpage
\section{Drawings}
\begin{figure}[H]
    \centering
    \includesvg[width=0.6\linewidth]{svg/arm_exp.svg}
    \caption{Exploded arm assembly}
    \label{fig:assembly}
\end{figure}

The arm is comprised of six assemblies. As can be seen in Figure \ref{fig:assembly}, this includes the base of the arm, four link pieces, and a plate at the end of the arm to visualise actuation of the final servo.

The base, shown in Figures \ref{fig:base} is designed to hold the arm's ESP32 microcontroller, power and servo control boards, and any weights that may want to be added in order to hold the arm in place during movement.

Figures \ref{fig:seg1}, \ref{fig:seg2}, \ref{fig:seg3}, and \ref{fig:seg4} show the construction of each link of the arm, and figure \ref{fig:seg5} shows the construction of the final "plate" link.

% --- Base ---
\begin{figure}[H]
    \centering
    \begin{subfigure}{0.48\linewidth}
        \centering
        \includesvg[width=0.8\linewidth]{svg/base.svg} 
        \subcaption{Base assembly}
        \label{fig:base_a}
    \end{subfigure}
    \hfill
    \begin{subfigure}{0.48\linewidth}
        \centering
        \includesvg[width=\linewidth]{svg/base_exp.svg} 
        \subcaption{Base exploded view}
        \label{fig:base_b}
    \end{subfigure}
    \caption{Base views}
    \label{fig:base}
\end{figure}

% --- Segment 1 ---
\begin{figure}[H]
    \centering
    \begin{subfigure}{0.48\linewidth}
        \centering
        \includesvg[width=0.4\linewidth]{svg/seg1.svg} 
        \subcaption{Segment 1 assembly}
        \label{fig:seg1_a}
    \end{subfigure}
    \hfill
    \begin{subfigure}{0.48\linewidth}
        \centering
        \includesvg[width=\linewidth]{svg/seg1_exp.svg} 
        \subcaption{Segment 1 exploded}
        \label{fig:seg1_b}
    \end{subfigure}
    \caption{Segment 1 views}
    \label{fig:seg1}
\end{figure}

% --- Segment 2 ---
\begin{figure}[H]
    \centering
    \begin{subfigure}{0.48\linewidth}
        \centering
        \includesvg[width=0.4\linewidth]{svg/seg2.svg} 
        \subcaption{Segment 2 assembly}
        \label{fig:seg2_a}
    \end{subfigure}
    \hfill
    \begin{subfigure}{0.48\linewidth}
        \centering
        \includesvg[width=\linewidth]{svg/seg2_exp.svg} 
        \subcaption{Segment 2 exploded}
        \label{fig:seg2_b}
    \end{subfigure}
    \caption{Segment 2 views}
    \label{fig:seg2}
\end{figure}

% --- Segment 3 ---
\begin{figure}[H]
    \centering
    \begin{subfigure}{0.48\linewidth}
        \centering
        \includesvg[width=0.4\linewidth]{svg/seg3.svg} 
        \subcaption{Segment 3 assembly}
        \label{fig:seg3_a}
    \end{subfigure}
    \hfill
    \begin{subfigure}{0.48\linewidth}
        \centering
        \includesvg[width=\linewidth]{svg/seg3_exp.svg} 
        \subcaption{Segment 3 exploded}
        \label{fig:seg3_b}
    \end{subfigure}
    \caption{Segment 3 views}
    \label{fig:seg3}
\end{figure}

% --- Segment 4 ---
\begin{figure}[H]
    \centering
    \begin{subfigure}{0.48\linewidth}
        \centering
        \includesvg[width=0.4\linewidth]{svg/seg4.svg} 
        \subcaption{Segment 4 assembly}
        \label{fig:seg4_a}
    \end{subfigure}
    \hfill
    \begin{subfigure}{0.48\linewidth}
        \centering
        \includesvg[width=\linewidth]{svg/seg4_exp.svg} 
        \subcaption{Segment 4 exploded}
        \label{fig:seg4_b}
    \end{subfigure}
    \caption{Segment 4 views}
    \label{fig:seg4}
\end{figure}

% --- Segment 5 ---
\begin{figure}[H]
    \centering
    \includesvg[width=0.2\linewidth]{svg/seg5.svg} 
    \caption{Segment 5 views}
    \label{fig:seg5}
\end{figure}


% --- 3: BILL OF MATERIALS -------------------------------------------------------------------------
\newpage
\section{Bill of Materials (BOM)}
\begin{table}[H]
    \centering
    \caption{Purchased Parts}
    \label{tab:bom_pp}
    \begin{tabularx}{\textwidth}{@{} l X r @{}}
        \toprule
        \textbf{Description} & \textbf{Qty} \\
        \midrule
        MG996R Servo                                              & 5 \\
        ESP32 Devkit Type-C                                       & 1 \\
        PCA9685 16-Channel 12-Bit PWM Driver                      & 1 \\
        DC5525 to Screw Terminal Connector DC Power Jack (Female) & 1 \\
        10A Adjustable DC-DC Step Down Converter                  & 1 \\
        \bottomrule
    \end{tabularx}
\end{table}

\begin{table}[H]
    \centering
    \caption{Manufactured Parts}
    \label{tab:bom_manufactured}
    \begin{tabularx}{\textwidth}{@{} l X r @{}}
        \toprule
        \textbf{Part} & \textbf{Material} & \textbf{Qty} \\
        \midrule
        \texttt{base}                   & PLA - White & 1 \\
        \texttt{seg1}                   & PLA - White & 1 \\
        \texttt{seg2a}                  & PLA - White & 1 \\
        \texttt{seg2b}                  & PLA - White & 1 \\
        \texttt{seg3a}                  & PLA - White & 1 \\
        \texttt{seg3b}                  & PLA - White & 1 \\
        \texttt{seg4a}                  & PLA - White & 1 \\
        \texttt{seg4b}                  & PLA - White & 1 \\
        \texttt{seg5}                   & PLA - White & 1 \\
        \texttt{seg1\_shell-side\_l}    & PLA - White & 1 \\
        \texttt{seg1\_shell-side\_r}    & PLA - White & 1 \\
        \texttt{seg2\_shell-side\_r}    & PLA - White & 1 \\
        \texttt{seg3\_shell-side\_l}    & PLA - White & 1 \\
        \texttt{seg4\_shell-side\_r}    & PLA - White & 1 \\
        \texttt{seg2\_strut}            & PLA - White & 2 \\
        \texttt{seg3\_strut}            & PLA - White & 1 \\
        \bottomrule
    \end{tabularx}
\end{table}

\begin{table}[H]
    \centering
    \caption{Standard Parts}
    \label{tab:bom_fasteners}
    \begin{tabularx}{\textwidth}{@{} l X r @{}}
        \toprule
        \textbf{Description} & \textbf{Qty} \\
        \midrule
        M3x6 Pan Head Screw  & 9  \\
        M3x10 Pan Head Screw & 22 \\
        M3x14 Pan Head Screw & 1  \\
        M3 Nut               & 3  \\
        \bottomrule
    \end{tabularx}
\end{table}


% --- 4: CALCULATIONS ------------------------------------------------------------------------------
% \newpage
% \section{Calculations}
% \subsection{Load Analysis}
% Using the dimensions from Figure \ref{fig:assembly}, the moment of inertia $I$ is calculated as:
% \[
%     I = \frac{bh^3}{12}
% \]
% Substituting the values $b = \SI{20}{\milli\meter}$ and $h = \SI{5}{\milli\meter}$:
% \[
%     I = \frac{20 \times 5^3}{12} = \SI{208.33}{\milli\meter\squared}
% \]

\end{document}